\documentclass[11pt,a4paper]{article}

\usepackage[utf8]{inputenc}
\usepackage[T1]{fontenc}
\usepackage{lmodern}
\usepackage[french]{babel}
\usepackage{geometry}
\usepackage{xcolor}
\usepackage{hyperref}
\usepackage{fancyhdr}
\usepackage{titlesec}
\usepackage{enumitem}
\usepackage{booktabs}
\usepackage{longtable}
\usepackage{array}
\usepackage{amsmath}
\usepackage{amssymb}
\usepackage{graphicx}
\usepackage{listings}
\usepackage{lastpage}

\geometry{margin=2.2cm}

\definecolor{accent}{HTML}{1F5F5B}
\definecolor{accentlight}{HTML}{E8F3F1}
\definecolor{codebg}{HTML}{F6F8FA}
\definecolor{codetext}{HTML}{24292F}
\definecolor{softgray}{HTML}{6B7280}
\definecolor{danger}{HTML}{8A1F11}

\hypersetup{
  colorlinks=true,
  linkcolor=accent,
  urlcolor=accent,
  citecolor=accent,
  pdftitle={Rapport de modification AC Video JEPA},
  pdfauthor={Codex}
}

\pagestyle{fancy}
\fancyhf{}
\lhead{\textcolor{softgray}{AC Video JEPA}}
\rhead{\textcolor{softgray}{Rapport technique}}
\cfoot{\textcolor{softgray}{Page \thepage\ sur \pageref{LastPage}}}
\renewcommand{\headrulewidth}{0.2pt}
\renewcommand{\footrulewidth}{0.2pt}

\titleformat{\section}
  {\Large\bfseries\color{accent}}
  {\thesection}{0.8em}{}
\titleformat{\subsection}
  {\large\bfseries\color{accent}}
  {\thesubsection}{0.8em}{}

\setlist[itemize]{leftmargin=1.2em, itemsep=0.25em}
\setlist[enumerate]{leftmargin=1.5em, itemsep=0.25em}

\lstdefinestyle{code}{
  backgroundcolor=\color{codebg},
  basicstyle=\ttfamily\small\color{codetext},
  breaklines=true,
  columns=fullflexible,
  keepspaces=true,
  frame=single,
  framerule=0pt,
  rulecolor=\color{codebg},
  xleftmargin=0.4em,
  xrightmargin=0.4em,
  aboveskip=0.7em,
  belowskip=0.7em
}
\lstset{style=code}

\newcommand{\filepath}[1]{\texttt{#1}}
\newcommand{\code}[1]{\texttt{#1}}

\newcommand{\infobox}[2]{
  \vspace{0.8em}
  \noindent
  \fcolorbox{accent}{accentlight}{
    \begin{minipage}{0.94\linewidth}
      \vspace{0.5em}
      \textbf{\color{accent}#1}\par\vspace{0.3em}
      #2
      \vspace{0.5em}
    \end{minipage}
  }
  \vspace{0.8em}
}

\begin{document}

\begin{titlepage}
  \centering
  \vspace*{2.0cm}
  {\Huge\bfseries\color{accent}Rapport technique\par}
  \vspace{0.4cm}
  {\LARGE Adaptation de V-JEPA-AC vers un prédicteur direct multi-horizon causal\par}
  \vspace{1.0cm}
  {\large Projet : \texttt{V-JEPA-AC-paral}\par}
  \vspace{0.2cm}
  {\large Date du rapport : \today\par}
  \vfill
  \infobox{Résumé}{
    Ce rapport décrit les modifications apportées au modèle \code{ac\_video\_jepa}.
    L'objectif est de remplacer le rollout autoregressif mono-step par une prédiction
    directe multi-horizon, calculée en un seul forward pass, tout en garantissant une
    causalité stricte sur la séquence d'actions.
  }
  \vfill
  {\small Généré dans le dossier \texttt{rapport/}.}
\end{titlepage}

\tableofcontents
\newpage

\section{Contexte et objectif}

Le code initial de l'exemple \filepath{examples/ac\_video\_jepa} utilisait un prédicteur
autoregressif de type GRU. Le fonctionnement était le suivant :

\[
z_t, a_t \rightarrow \hat z_{t+1}
\]

puis :

\[
\hat z_{t+1}, a_{t+1} \rightarrow \hat z_{t+2}
\]

et ainsi de suite. Ce choix est simple, mais il introduit une accumulation d'erreur :
une erreur sur \(\hat z_{t+1}\) devient l'état d'entrée de la prédiction suivante.

\subsection{Nouvel objectif}

L'objectif demandé était de transformer le modèle vers une formulation directe
multi-horizon :

\[
G(z_t, s_t, a_t, \dots, a_{t+K-1})
\rightarrow
(\hat z_{t+1}, \dots, \hat z_{t+K})
\]

avec \(K=5\) dans la configuration par défaut.

La dépendance voulue est :

\begin{center}
\begin{tabular}{ll}
\toprule
Sortie & Dépendances autorisées \\
\midrule
\(\hat z_{t+1}\) & \(z_t, s_t, a_t\) \\
\(\hat z_{t+2}\) & \(z_t, s_t, a_t, a_{t+1}\) \\
\(\hat z_{t+3}\) & \(z_t, s_t, a_t, a_{t+1}, a_{t+2}\) \\
\(\hat z_{t+4}\) & \(z_t, s_t, a_t, a_{t+1}, a_{t+2}, a_{t+3}\) \\
\(\hat z_{t+5}\) & \(z_t, s_t, a_t, a_{t+1}, a_{t+2}, a_{t+3}, a_{t+4}\) \\
\bottomrule
\end{tabular}
\end{center}

Le point essentiel est que toutes les sorties sont produites en un seul forward pass,
mais chaque horizon ne voit que le préfixe d'actions qui lui correspond.

\section{Vue d'ensemble des modifications}

Huit fichiers ont été modifiés.

\begin{center}
\begin{tabular}{>{\small\raggedright\arraybackslash}p{0.54\linewidth}p{0.34\linewidth}}
\toprule
Fichier & Rôle de la modification \\
\midrule
\filepath{eb\_jepa/architectures.py} & Ajout du masque causal et du nouveau prédicteur Transformer multi-horizon. \\
\filepath{eb\_jepa/jepa.py} & Ajout du mode \code{direct\_multi\_horizon} dans \code{JEPA.unroll}. \\
\filepath{eb\_jepa/losses.py} & Ajout de la loss pondérée \code{MultiHorizonLoss}. \\
\filepath{eb\_jepa/planning.py} & Adaptation du planner pour utiliser automatiquement le nouveau mode et ajout de régularisations d'actions. \\
\filepath{examples/ac\_video\_jepa/main.py} & Branchement du nouveau prédicteur dans l'entraînement AC Video JEPA. \\
\filepath{examples/ac\_video\_jepa/cfgs/train.yaml} & Configuration par défaut vers \(K=5\) et prédicteur multi-horizon. \\
\filepath{examples/ac\_video\_jepa/cfgs/planning\_mppi.yaml} & MPC aligné sur horizon 5, coût terminal, régularisation des actions. \\
\filepath{tests/test\_jepa\_output\_formats.py} & Ajout de tests de masque causal et de sortie multi-horizon. \\
\bottomrule
\end{tabular}
\end{center}

\section{Nouveau prédicteur causal multi-horizon}

\subsection{Fichier modifié}

\filepath{eb\_jepa/architectures.py}

\subsection{Masque causal}

Une fonction a été ajoutée :

\begin{lstlisting}[language=Python]
build_multi_horizon_mask(
    num_context_tokens,
    num_patches,
    horizon,
    device,
)
\end{lstlisting}

Elle construit un masque d'attention booléen. Dans cette convention, \code{True}
signifie que l'attention est bloquée.

L'ordre des tokens est :

\[
[\text{tokens contexte latent}]
+
[\text{tokens actions}]
+
[\text{tokens queries futures}]
\]

Le masque impose trois règles.

\begin{enumerate}
  \item Les tokens de contexte latent ne voient que le contexte latent.
  \item Le token d'action \(a_h\) voit le contexte et le préfixe \(a_0,\dots,a_h\).
  \item La query \(q_h\), qui prédit \(\hat z_{t+h+1}\), voit le contexte et \(a_0,\dots,a_h\), mais pas les actions futures.
\end{enumerate}

\infobox{Pourquoi ce masque est nécessaire}{
  Il ne suffit pas de bloquer directement l'accès de \(q_1\) à \(a_4\). Il faut aussi
  empêcher les tokens intermédiaires de transporter indirectement l'information de
  \(a_4\). C'est pour cela que les tokens d'action eux-mêmes sont causaux.
}

\subsection{Classe \code{CausalMultiHorizonPredictor}}

Une nouvelle classe a été ajoutée :

\begin{lstlisting}[language=Python]
class CausalMultiHorizonPredictor(nn.Module):
    ...
\end{lstlisting}

Son interface interne accepte :

\begin{lstlisting}[language=Python]
state:  [B, D, context, H, W]
action: [B, A, K]
return: [B, D, K, H, W]
\end{lstlisting}

Dans le chemin AC direct utilisé ici, si l'encodeur produit une carte spatiale
\([B,D,T,H,W]\), cette carte est d'abord moyennée spatialement :

\[
[B,D,T,H,W] \rightarrow [B,D,T,1,1]
\]

Le prédicteur reçoit donc un token global par image. Dans le cas actuel de l'encodeur
Impala, la représentation est déjà globale :

\[
z_t \in \mathbb{R}^{512 \times 1 \times 1}
\]

Donc le prédicteur travaille principalement avec :

\[
[B, 512, context, 1, 1] + [B, 2, K]
\rightarrow
[B, 512, K, 1, 1]
\]

Une variante dense par patches resterait possible en configurant le prédicteur avec
plusieurs tokens spatiaux, mais ce n'est pas le comportement par défaut de ce projet.

\section{Nouveau mode dans \code{JEPA.unroll}}

\subsection{Fichier modifié}

\filepath{eb\_jepa/jepa.py}

\subsection{Mode ajouté}

Un nouveau mode a été ajouté :

\begin{lstlisting}[language=Python]
unroll_mode="direct_multi_horizon"
\end{lstlisting}

Ce mode conserve les anciens modes :

\begin{itemize}
  \item \code{parallel}, utilisé par l'exemple Video JEPA non actionné ;
  \item \code{autoregressive}, utilisé par l'ancien AC Video JEPA GRU ;
  \item \code{direct\_multi\_horizon}, ajouté pour le nouvel objectif.
\end{itemize}

\subsection{Comportement en entraînement}

Pendant l'entraînement, le code encode d'abord toute la séquence :

\[
z_{0:T} = E(x_{0:T})
\]

Puis il crée des fenêtres glissantes :

\[
z_{t-c+1:t},\quad a_t,\dots,a_{t+K-1}
\rightarrow
z_{t+1},\dots,z_{t+K}
\]

où \(c\) est \code{context\_length}.

Avec la configuration actuelle :

\[
c = 1,\quad K = 5
\]

ce qui donne :

\[
z_t,\quad a_t,\dots,a_{t+4}
\rightarrow
z_{t+1},\dots,z_{t+5}
\]

L'intérêt est que chaque horizon est appris directement depuis le vrai latent courant,
et non depuis une prédiction précédente.

\subsection{Comportement en planning}

En planning, le mode reçoit :

\[
z_t,\quad a_t,\dots,a_{t+K-1}
\]

et produit :

\[
\hat z_{t+1:t+K}
\]

Si le planner demande un horizon plus long que celui du prédicteur, le code découpe
en blocs. Pour la configuration par défaut, le planning utilise exactement \(K=5\).

\section{Nouvelle loss multi-horizon}

\subsection{Fichier modifié}

\filepath{eb\_jepa/losses.py}

\subsection{Classe ajoutée}

\begin{lstlisting}[language=Python]
class MultiHorizonLoss(nn.Module):
    ...
\end{lstlisting}

Elle reçoit :

\begin{lstlisting}[language=Python]
pred:   [B, D, K, H, W]
target: [B, D, K, H, W]
\end{lstlisting}

Elle applique une loss par horizon, puis moyenne les horizons avec des poids uniformes
par défaut.

Par défaut :

\[
\gamma = 1.0
\]

Pour \(K=5\), les poids normalisés sont équivalents à :

\[
[0.2,\;0.2,\;0.2,\;0.2,\;0.2]
\]

\subsection{Raison du choix}

La pondération uniforme force le modèle à apprendre une dynamique cohérente à chaque
horizon, plutôt qu'à privilégier uniquement le latent terminal. Cela réduit le risque
d'apprendre un raccourci entre une séquence d'actions et un état final plausible.

\section{Branchement dans AC Video JEPA}

\subsection{Fichier modifié}

\filepath{examples/ac\_video\_jepa/main.py}

\subsection{Nouveau choix de prédicteur}

Le script d'entraînement lit maintenant :

\begin{lstlisting}[language=Python]
predictor_type = cfg.model.get("predictor_type", "direct_multi_horizon")
\end{lstlisting}

Deux chemins existent :

\begin{itemize}
  \item \code{predictor\_type: rnn} conserve l'ancien \code{RNNPredictor}.
  \item \code{predictor\_type: direct\_multi\_horizon} active le nouveau prédicteur.
\end{itemize}

Cela permet de comparer l'ancien modèle et le nouveau sans supprimer le comportement
précédent.

\subsection{Loss utilisée}

Si le nouveau mode est actif, le training utilise :

\begin{lstlisting}[language=Python]
MultiHorizonLoss(...)
\end{lstlisting}

Sinon, il conserve :

\begin{lstlisting}[language=Python]
SquareLossSeq()
\end{lstlisting}

\section{Modifications de configuration}

\subsection{Training}

Fichier :

\filepath{examples/ac\_video\_jepa/cfgs/train.yaml}

Les changements principaux sont :

\begin{lstlisting}
nsteps: 5
predictor_type: direct_multi_horizon
context_length: 1
pred_dim: 512
predictor_depth: 4
predictor_heads: 8
horizon_loss_gamma: 1.0
horizon_loss_type: smooth_l1
\end{lstlisting}

La valeur \(K=5\) correspond directement à la formulation discutée.

\subsection{Planning MPPI}

Fichier :

\filepath{examples/ac\_video\_jepa/cfgs/planning\_mppi.yaml}

Les changements principaux sont :

\begin{lstlisting}
plan_length: 5
sum_all_diffs: false
action_l2_coeff: 0.01
action_smoothness_coeff: 0.05
\end{lstlisting}

Le coût MPC devient principalement terminal :

\[
J = d(\hat z_{t+5}, z_{\text{goal}})
\]

avec régularisation légère des actions.

\section{Adaptation du planner}

\subsection{Fichier modifié}

\filepath{eb\_jepa/planning.py}

\subsection{Détection automatique du nouveau prédicteur}

Le planner détecte :

\begin{lstlisting}[language=Python]
predictor.direct_multi_horizon
\end{lstlisting}

Si ce flag existe, il appelle :

\begin{lstlisting}[language=Python]
unroll_mode="direct_multi_horizon"
\end{lstlisting}

Sinon, il conserve l'ancien mode :

\begin{lstlisting}[language=Python]
unroll_mode="autoregressive"
\end{lstlisting}

\subsection{Régularisation des actions}

Deux coefficients optionnels ont été ajoutés :

\begin{lstlisting}[language=Python]
action_l2_coeff
action_smoothness_coeff
\end{lstlisting}

Le coût devient :

\[
J =
d(\hat z_{t+K}, z_g)
+
\lambda \frac{1}{K}\sum_i \|a_i\|^2
+
\mu \frac{1}{K-1}\sum_i \|a_i-a_{i-1}\|^2
\]

Cela stabilise le planning et évite des séquences d'actions inutilement violentes.

\section{Tests ajoutés}

\subsection{Fichier modifié}

\filepath{tests/test\_jepa\_output\_formats.py}

\subsection{Test du masque}

Test ajouté :

\begin{lstlisting}[language=Python]
test_direct_multi_horizon_mask()
\end{lstlisting}

Il vérifie que :

\begin{itemize}
  \item les actions ne voient pas les actions futures ;
  \item les queries ne voient pas les actions futures ;
  \item les queries ne communiquent pas entre elles, sauf auto-attention à elles-mêmes.
\end{itemize}

\subsection{Test du nouveau \code{unroll}}

Test ajouté :

\begin{lstlisting}[language=Python]
test_direct_multi_horizon_unroll_with_loss()
\end{lstlisting}

Il vérifie :

\begin{itemize}
  \item la forme de sortie en entraînement ;
  \item la présence d'une loss finie ;
  \item la forme de sortie en planning sans loss.
\end{itemize}

\section{Validation effectuée}

Les validations suivantes ont été exécutées avec succès :

\begin{lstlisting}
python3 -m py_compile ...
git diff --check
\end{lstlisting}

Cela garantit :

\begin{itemize}
  \item absence d'erreurs syntaxiques Python ;
  \item absence de problèmes de whitespace dans le diff.
\end{itemize}

\subsection{Limite de validation}

Les tests PyTorch complets n'ont pas pu être lancés dans l'environnement courant, car :

\begin{itemize}
  \item \code{torch} n'est pas installé dans le \code{python3} disponible ;
  \item \code{pytest} n'est pas disponible ;
  \item \code{uv} n'est pas disponible ;
  \item le \code{python3} système est en version 3.9.6, alors que le projet demande Python 3.12.
\end{itemize}

\section{Limites et points à surveiller}

\subsection{Historique réel en MPC}

Le code supporte \code{context\_length > 1} au niveau du prédicteur et du training.
Cependant, le planner actuel reçoit principalement l'observation courante.

Si l'objectif final est vraiment :

\[
G(z_{t-c+1:t}, a_t,\dots,a_{t+K-1})
\]

alors il faudra ajouter une mémoire d'observations dans \code{GCAgent}, afin que le MPC
fournisse les derniers états réels au modèle au lieu d'un contexte paddé.

\subsection{Pooling spatial avant prédiction}

L'encodeur Impala actuel produit déjà une représentation globale :

\[
[B, 512, T, 1, 1]
\]

Si un encodeur V-JEPA/V-JEPA2 dense est branché plus tard et produit une grille de tokens
spatiaux, le chemin \code{direct\_multi\_horizon} applique maintenant un mean pooling
sur \(H,W\) avant le prédicteur et avant la loss de prédiction.

\subsection{Comparaison expérimentale à faire}

Il faudra comparer :

\begin{itemize}
  \item ancien \code{RNNPredictor} autoregressif ;
  \item nouveau prédicteur direct multi-horizon ;
  \item \code{context\_length = 1} contre \code{context\_length > 1} ;
  \item coût MPC terminal seul contre coût sur tous les horizons.
\end{itemize}

\section{Commandes utiles}

Pour lancer l'entraînement avec le nouveau mode :

\begin{lstlisting}
python -m examples.ac_video_jepa.main \
  --fname examples/ac_video_jepa/cfgs/train.yaml
\end{lstlisting}

Pour revenir à l'ancien prédicteur GRU :

\begin{lstlisting}
python -m examples.ac_video_jepa.main \
  --fname examples/ac_video_jepa/cfgs/train.yaml \
  model.predictor_type=rnn
\end{lstlisting}

Pour tester l'idée avec un contexte plus long :

\begin{lstlisting}
python -m examples.ac_video_jepa.main \
  --fname examples/ac_video_jepa/cfgs/train.yaml \
  model.context_length=3
\end{lstlisting}

\section{Conclusion}

La modification principale est conceptuelle : le modèle n'est plus obligé de propager
ses propres erreurs horizon après horizon. Il apprend maintenant une fonction directe :

\[
G(z_t, a_t,\dots,a_{t+4})
\rightarrow
(\hat z_{t+1},\dots,\hat z_{t+5})
\]

avec une causalité stricte :

\[
\hat z_{t+h}
\not\leftarrow
a_{t+h}, a_{t+h+1}, \dots, a_{t+4}
\]

Ce changement aligne mieux le modèle avec le MPC : le planner peut scorer directement
une séquence candidate d'actions via l'état latent terminal prédit, puis exécuter seulement
la première action et replanifier.

\end{document}
